\chapter{Аналіз прогалин у навичках}

\section{Методика оцінки}

Для оцінки покриття AI та цифрових навичок розроблено фреймворк з 11
категорій навичок. Покриття оцінювалось за наявністю ключових слів
у назвах курсів та текстах силабусів.

\section{Результати аналізу}

\begin{figure}[H]
\centering
\includegraphics[width=0.9\textwidth]{figures/gap_analysis.png}
\caption{Покриття AI та цифрових навичок}
\label{fig:gap_analysis}
\end{figure}

============================================================
AI \& DIGITAL SKILLS GAP ANALYSIS
============================================================

\#\# COVERAGE SUMMARY

Skills checked: 10
Fully covered: 7
Partially covered: 0
Missing: 3

\#\# COVERAGE SCORES

AI Core Skills: 0.0\%
Digital Tools: 71.4\%
Overall: 35.7\%

\#\# MISSING SKILLS

  ✗ AI Fundamentals
  ✗ Generative AI \& Design
  ✗ Computer Vision

\#\# SKILL DETAILS

✓ 3D Design \& Modeling: 100.0\%
    Courses: 20 Технології 3D-моделювання, 3D-моделювання, Інженерно-будівельна та комп’ютерна графіка (ArchiCAD, AutoCAD)
✓ Motion Design \& Animation: 100.0\%
    Courses: Фото та відео монтаж, Основи фотовідеодизайну, і анімація проєкту
✓ UX/UI Design: 100.0\%
    Courses: 24 Дизайн інтерфейсів, Дизайн інтерфейсів, робота «Веб-технології та веб-дизайн»
✓ Interactive \& Computational Design: 100.0\%
    Courses: Web\_дизайн\_та\_web\_програмування, Шовкодрук, 23 Web дизайн та web програмування
✓ 2D Digital Design Tools: 33.3\%
    Courses: Цифрова ілюстрація і анімація
✓ Data Visualization: 33.3\%
    Courses: Інфографіка
✓ VR/AR Design: 33.3\%
    Courses: English for career development, Інженерно-будівельна та комп’ютерна графіка (ArchiCAD, AutoCAD)
✗ AI Fundamentals: 0.0\%
✗ Generative AI \& Design: 0.0\%
✗ Computer Vision: 0.0\%

\#\# RECOMMENDATIONS

1. Add introductory course on AI/ML basics for designers, covering concepts without deep technical requirements
2. Integrate generative AI tools (Midjourney, DALL-E, Stable Diffusion) into existing design courses or create dedicated workshop

============================================================

\section{Виявлені прогалини}

Найбільші прогалини виявлено в таких областях:
\begin{itemize}
    \item Генеративний AI та prompt engineering
    \item Параметричний та інтерактивний дизайн
    \item VR/AR дизайн
    \item Візуалізація даних
\end{itemize}

\section{Порівняння за рівнями}

Магістерські програми демонструють дещо вище покриття цифрових навичок,
проте AI-інструменти залишаються недостатньо представленими на всіх рівнях.
